\documentclass[../../main.tex]{subfiles}% 注意这里的文件路径不能用 ./main.tex ,否则用latexmk编译子文件会报错
\graphicspath{{\subfix{./image/}}{\subfix{../../image/}}} % 指定图片引用路径,后续可以直接使用图片文件名
% 如果图片存放在上级目录,则使用 ../../image/ 作为备用路径

% 例如:
% \begin{figure}[H]
% \centering
% \includegraphics[width=0.3\textwidth]{图.png}
% \caption{}
% \label{figure:图}
% \end{figure}
% 注意:上述\label{}一定要放在\caption{}之后,否则引用图片序号会只会显示??.

\begin{document}

\section{平面上的Dirichlet问题}

\begin{definition}
设 \(u\) 是圆周 \(|z|=R\) 上的逐段连续函数，称
\begin{align}
P[u](z) = \int_{0}^{2\pi} P(z, Re^{i\theta})u(Re^{i\theta}) \mathrm{d}\theta. \label{eq:poisson_def}
\end{align}
为 \(u\) 在圆盘 \(B(0,R)\) 中的\textbf{Poisson 积分}.
\end{definition}

\begin{theorem}\label{theorem:poisson_solution}
设 \(u\) 是圆周 \(|z|=R\) 上的逐段连续函数，那么 \(P[u](z)\) 是圆盘 \(B(0,R)\) 中的有界调和函数，且在 \(u\) 的连续点 \(Re^{i\varphi_0}\) 处有
\begin{align}
\lim\limits_{z\to Re^{i\varphi_0}} P[u](z) = u(Re^{i\varphi_0}). \label{eq:thm1_boundary}
\end{align}
\end{theorem}
\begin{proof}
由\subcref{proposition:poisson_kernel}{proposition:poisson_kernel-3}即知 \(P[u](z)\) 是 \(B(0,R)\) 中的调和函数. 因为 \(u\) 有界，设 \(|u(Re^{i\varphi})|\leqslant M, 0\leqslant \varphi \leqslant 2\pi.\) 由\subcref{proposition:poisson_kernel}{proposition:poisson_kernel-1}和\subcref{proposition:poisson_kernel}{proposition:poisson_kernel-2}即得
\begin{align*}
|P[u](z)| \leqslant \int_{0}^{2\pi} P(z, Re^{i\varphi}) |u(Re^{i\varphi})| \mathrm{d}\varphi \leqslant M.
\end{align*}
这就证明了 \(P[u](z)\) 是 \(B(0,R)\) 中的有界调和函数.
下面证明等式 \eqref{eq:thm1_boundary}. 因为 \(u\) 在 \(Re^{i\varphi_0}\) 处连续，故对任意 \(\varepsilon>0\)，存在 \(\delta>0\)，当 \(|\varphi-\varphi_0|\leqslant \delta\) 时，有
\begin{align}
|u(Re^{i\varphi}) - u(Re^{i\varphi_0})| < \varepsilon. \label{eq:continuity}
\end{align}
由\subcref{proposition:poisson_kernel}{proposition:poisson_kernel-2}，可得
\begin{align*}
u(Re^{i\varphi_0}) = \int_{0}^{2\pi} P(z, Re^{i\varphi}) u(Re^{i\varphi_0}) \mathrm{d}\varphi.
\end{align*}
因而
\begin{align*}
|P[u](z) - u(Re^{i\varphi_0})| &\leqslant \int_{0}^{2\pi} P(z, Re^{i\varphi}) |u(Re^{i\varphi}) - u(Re^{i\varphi_0})| \mathrm{d}\varphi \\
&= \int_{|\varphi-\varphi_0|\leqslant \delta} P(z, Re^{i\varphi}) |u(Re^{i\varphi}) - u(Re^{i\varphi_0})| \mathrm{d}\varphi \\
&\quad + \int_{|\varphi-\varphi_0|>\delta} P(z, Re^{i\varphi}) |u(Re^{i\varphi}) - u(Re^{i\varphi_0})| \mathrm{d}\varphi \\
&= I_1 + I_2.
\end{align*}
由 \eqref{eq:continuity} 式立刻可得
\begin{align*}
I_1 &= \int_{|\varphi-\varphi_0|\leqslant \delta} P(z, Re^{i\varphi}) |u(Re^{i\varphi}) - u(Re^{i\varphi_0})| \mathrm{d}\varphi \\
&< \varepsilon \int_{0}^{2\pi} P(z, Re^{i\varphi}) \mathrm{d}\varphi = \varepsilon.
\end{align*}
若令 \(z=re^{i\theta}\)，则
\begin{align*}
P(z, Re^{i\varphi}) = \frac{1}{2\pi} \frac{R^2 - r^2}{R^2 - 2rR\cos(\theta-\varphi) + r^2}.
\end{align*}
当 \(z=re^{i\theta} \to Re^{i\varphi_0}\) 时，由于 \(r \to R, \theta \to \varphi_0\)，因而 \(R^2 - r^2 \to 0\)，\(R^2 - 2rR\cos(\theta-\varphi) + r^2 \to 2R^2(1 - \cos(\varphi_0-\varphi))\).
而当 \(|\varphi-\varphi_0|>\delta\) 时，\(2R^2(1 - \cos(\varphi_0-\varphi)) > 2R^2(1 - \cos\delta)\). 故对任意的 \(\varepsilon>0\) 及固定的 \(\delta>0\)，必存在 \(\eta>0\)，使得当 \(|z-Re^{i\varphi_0}|<\eta\) 且 \(|\varphi-\varphi_0|>\delta\) 时，有
\begin{gather*}
R^2 - 2rR\cos(\theta-\varphi) + r^2 > 2R^2(1-\cos\delta), \\
R^2 - r^2 < \frac{R^2}{M}(1-\cos\delta)\varepsilon.
\end{gather*}
于是
\begin{align*}
I_2 &= \frac{1}{2\pi} \int_{|\varphi-\varphi_0|>\delta} \frac{R^2 - r^2}{R^2 - 2rR\cos(\theta-\varphi) + r^2} |u(Re^{i\varphi}) - u(Re^{i\varphi_0})| \mathrm{d}\varphi \\
&< \frac{1}{2\pi} \cdot 2M \cdot \frac{R^2}{M} \cdot \frac{(1-\cos\delta)\varepsilon}{2R^2(1-\cos\delta)} \cdot 2\pi = \varepsilon.
\end{align*}
由上述两式即知 \eqref{eq:thm1_boundary} 式成立.
\end{proof}

\begin{definition}\label{definition:dirichlet_disk}
设 \(u\) 是圆周 \(|z|=R\) 上的有界逐段连续函数. 圆盘 \(B(0,R)\) 上的 Dirichlet 问题是指：在 \(B(0,R)\) 内寻找一个调和函数 \(h\)，使得该函数在闭圆盘 \(\overline{B(0,R)}\) 上连续，并且满足边界条件
\begin{align*}
\lim\limits_{\substack{z\to \zeta \\ z\in B(0,R)}} h(z) = u(\zeta), \quad \forall \zeta\in \{\zeta: |\zeta|=R\}.
\end{align*}
\end{definition}

\begin{proposition}\label{proposition:bounded_solution}
设 \(u\) 是圆周 \(|z|=R\) 上的逐段连续函数，\(E\) 是 \(u\) 的全体连续点的集合. 如果
\begin{align*}
M = \sup\limits_{\zeta\in E} u(\zeta), \quad m = \inf\limits_{\zeta\in E} u(\zeta),
\end{align*}
那么 \(u\) 在 \(B(0,R)\) 中 Dirichlet 问题的有界解 \(h\) 满足
\begin{align*}
m \leqslant h(z) \leqslant M, \quad z\in B(0,R).
\end{align*}
\end{proposition}
\begin{proof}
设 \(u\) 在 \(|z|=R\) 上的第一类间断点为 \(\zeta_1, \cdots, \zeta_n\). 令
\begin{align*}
v(z) = M + \varepsilon \sum\limits_{k=1}^{n} \ln \frac{2R}{|z-\zeta_k|},
\end{align*}
这里 \(\varepsilon\) 是任意一个正数. 容易知道 \(v(z)\) 是 \(B(0,R)\) 中的调和函数，且对任意 \(z\in B(0,R)\)，有 \(v(z)>M\). 取充分小的 \(r>0\)，令 \(D_k = B(\zeta_k,r)\cap B(0,R)\)，记
\begin{align*}
G_r = B(0,R) \setminus \bigcup\limits_{k=1}^{n} D_k.
\end{align*}
显然 \(v(z)-h(z)\) 是 \(G_r\) 中的调和函数，而且在 \(\overline{G_r}\) 上连续. 当 \(\zeta\in \partial G_r \cap \partial B(0,R)\) 时，有
\begin{align*}
v(\zeta) - h(\zeta) \geqslant M - h(\zeta) = M - u(\zeta) \geqslant 0.
\end{align*}
当 \(\zeta\in \partial G_r \cap B(0,R)\) 时，令 \(r\to 0\)，则 \(\zeta\to \zeta_k\)，\(v(\zeta)\to \infty\)，而 \(h\) 是一个有界函数，因而也有 \(v(\zeta)-h(\zeta)>0\). 总之，当 \(\zeta\in \partial G_r\) 时，\(v(\zeta)-h(\zeta)\geqslant 0\). 于是，由调和函数的极值原理，当 \(z\in G_r\) 时，有 \(v(z)-h(z)\geqslant 0\). 因为 \(r\) 可以任意小，所以这个不等式对任意 \(z\in B(0,R)\) 成立. 固定 \(z\)，让 \(\varepsilon\to 0\)，即得 \(h(z)\leqslant M\). 因为 \(-h\) 也是调和函数，类似地可以证得 \(-h(z)\leqslant -m\)，因而 \(m\leqslant h(z)\leqslant M\).
\end{proof}

\begin{theorem}
设 \(u\) 是圆周 \(|z|=R\) 上的逐段连续函数，那么 \(u\) 的 Poisson 积分 \(P[u](z)\) 是 Dirichlet 问题的唯一解.
\end{theorem}
\begin{proof}
\(P[u](z)\) 是 Dirichlet 问题的解已由 \cref{theorem:poisson_solution} 证明. 如果另外还有一个解 \(h\)，那么 \(P[u](z)-h(z)\) 是 \(B(0,R)\) 中的有界调和函数，且在 \(u\) 的连续点处取零值. 由 \cref{proposition:bounded_solution} 知，在 \(B(0,R)\) 中有 \(P[u](z)\equiv h(z)\). 这就证明了解的唯一性.
\end{proof}

\begin{theorem}
域 \(D\) 上具有平均值性质的函数一定是调和函数.
\end{theorem}
\begin{proof}
设 \(u\) 是域 \(D\) 上具有平均值性质的函数，因而对任意 \(B(z_0,r)\subset D\)，平均值性质成立. 根据 \cref{theorem:poisson_solution}，存在 \(B(z_0,r)\) 中的调和函数 \(v\)，它在 \(B(z_0,r)\) 上连续，在圆周 \(|z-z_0|=r\) 上与 \(u\) 相等. 由于 \(u,v\) 在 \(B(z_0,r)\) 中都有平均值性质，因而 \(h=u-v\) 也有平均值性质. 由\cref{proposition:mean_prop_extremum}，极值原理对 \(h\) 成立. 由于 \(h\) 在圆周 \(|z-z_0|=r\) 上恒等于零，因而在 \(B(z_0,r)\) 中 \(u(z)\equiv v(z)\)，故 \(u\) 是 \(B(z_0,r)\) 中的调和函数. 由于 \(B(z_0,r)\) 是包含在 \(D\) 中的任意圆盘，故 \(u\) 是 \(D\) 中的调和函数.
\end{proof}



\end{document}